\chapter{多元随机变量及其分布}\label{chap:multivariate-distributions}

\section*{Chapter 2 概述}

\textbf{本章路线图}：联合分布 $\to$ 边缘分布 $\to$ 条件分布 $\to$ 独立性 $\to$ 相关性 $\to$ 多元正态

\section{2.1 从一元到多元：为什么需要联合分布？}\label{sec:2.1}

在一元情形，我们用 cdf $F_X(x) = \mathbb{P}(X \leq x)$ 描述单个随机变量的行为。但在许多场景中，\emph{单个变量的分布并不能回答我们的问题}——我们需要知道多个变量如何共同变化。

\textbf{核心问题}：给定两个随机变量 $X_1$ 和 $X_2$，如何描述它们的联合行为？

这促使我们建立联合分布理论。本章以二元情形为主，揭示多元分布的核心思想。

\subsection{随机向量与联合分布}

\begin{definition}[随机向量]\label{def:random-vector-ch2}
设样本空间为 $\mathcal{C}$，$X_1$ 和 $X_2$ 为定义在 $\mathcal{C}$ 上的两个随机变量。则 $(X_1, X_2)$ 称为随机向量，其取值空间为 $\mathcal{D} = \{(x_1, x_2) : x_1 = X_1(c), x_2 = X_2(c), c \in \mathcal{C}\}$。
\end{definition}

\textbf{联合 cdf} 是一元 cdf 的自然推广：
\begin{equation}\label{eq:joint-cdf-ch2}
F_{X_1, X_2}(x_1, x_2) = \mathbb{P}(X_1 \leq x_1, X_2 \leq x_2).
\end{equation}

\textbf{离散型} 用 pmf 描述：
\begin{equation}\label{eq:joint-pmf-ch2}
p_{X_1, X_2}(x_1, x_2) = \mathbb{P}(X_1 = x_1, X_2 = x_2),
\end{equation}
满足 $p \geq 0$ 且 $\sum\sum_{\mathcal{D}} p = 1$。

\textbf{连续型} 用 pdf 描述：
\begin{equation}\label{eq:joint-pdf-ch2}
F_{X_1, X_2}(x_1, x_2) = \int_{-\infty}^{x_1} \int_{-\infty}^{x_2} f(w_1, w_2) \, dw_1 \, dw_2,
\end{equation}
其中 $f \geq 0$ 且 $\int\int_{\mathcal{D}} f = 1$。

\begin remark}\label{rmk:pdf-can-be-greater-than-1}
与一元情形不同，这里的 $f(x_1, x_2)$ 可以大于 1！因为它是"密度"而非概率值。
\end remark}

几何直观：若把联合 pdf $z = f(x_1, x_2)$ 视作一张曲面，则 $\mathbb{P}((X_1, X_2) \in A)$ 就是该曲面下区域 $A$ 上的体积。这是一元情形"曲线下面积"在二维的推广。

\subsection{边缘分布：积分即求和}\label{sec:marginal-distribution}

\textbf{核心问题}：已知 $(X_1, X_2)$ 的联合分布，如何得到单个 $X_1$ 的分布？

直觉上，我们需要"忽略" $X_2$ 的信息。具体做法是：
\begin{itemize}
\item \textbf{离散型}：对 $x_2$ 求和
\item \textbf{连续型}：对 $x_2$ 积分
\end{itemize}

这引导我们得到边缘分布公式：\footnote{推导见附录 \cref{sec:derivation-marginal-distributions}。}

\textbf{边缘 pmf}：
\begin{equation}\label{eq:marginal-pmf-ch2}
p_{X_1}(x_1) = \sum_{x_2} p_{X_1, X_2}(x_1, x_2).
\end{equation}

\textbf{边缘 pdf}：
\begin{equation}\label{eq:marginal-pdf-ch2}
f_{X_1}(x_1) = \int_{-\infty}^{\infty} f_{X_1, X_2}(x_1, x_2) \, dx_2.
\end{equation}

\begin remark}[为什么叫"边缘"？]
在二维离散情形，将联合 pmf 列表，边缘分布恰好是每行（或每列）的"边缘"之和。这解释了名称的由来。
\end remark}

\subsection{期望的推广与线性性}\label{sec:expectation-linear}

设 $Y = g(X_1, X_2)$ 是新的随机变量，其期望为
\[
\mathbb{E}Y = \int\int g(x_1, x_2) f(x_1, x_2) \, dx_1 \, dx_2 \quad \text{（连续型）}.
\]

关键发现：期望算子 $E$ 对加法保持线性，无论 $X_1, X_2$ 是否独立：\footnote{推导见附录 \cref{sec:derivation-expectation-linear}。}

\begin{equation}\label{th:expectation-linear-ch2}
\mathbb{E}(aY_1 + bY_2) = a\mathbb{E}Y_1 + b\mathbb{E}Y_2.
\end{equation}

\textbf{这为什么重要？} 在第1章我们学过独立随机变量之和的方差公式 $\text{var}(X_1 + X_2) = \text{var}(X_1) + \text{var}(X_2)$。期望的线性性是其基础——它无需独立性即可成立。

\section{2.2 变换：寻找新视角}\label{sec:2.2}

给定 $(X_1, X_2)$ 的分布，我们经常想研究其函数 $Y = u(X_1, X_2)$ 的分布。例如：
\begin{itemize}
\item 和 $Y = X_1 + X_2$（两个随机变量之和）
\item 商 $Y = X_1 / X_2$（比值）
\item 极坐标变换 $(R, \Theta) = (\sqrt{X_1^2 + X_2^2}, \arctan(X_2/X_1))$
\end{itemize}

对于连续型，一一对应的变换有简洁的公式：\footnote{推导见附录 \cref{sec:derivation-jacobian}。}

\begin{equation}\label{eq:transformation-joint-pdf-ch2}
f_{Y_1, Y_2}(y_1, y_2) = f_{X_1, X_2}(v_1(y_1, y_2), v_2(y_1, y_2)) \cdot |J|.
\end{equation}

其中 Jacobian 行列式
\begin{equation}\label{eq:jacobian-ch2}
J = \det\begin{pmatrix}
\frac{\partial x_1}{\partial y_1} & \frac{\partial x_1}{\partial y_2} \\
\frac{\partial x_2}{\partial y_1} & \frac{\partial x_2}{\partial y_2}
\end{pmatrix}
= \frac{\partial x_1}{\partial y_1} \cdot \frac{\partial x_2}{\partial y_2} - \frac{\partial x_1}{\partial y_2} \cdot \frac{\partial x_2}{\partial y_1}.
\end{equation}

\textbf{为什么有 $|J|$？} 这是面积元素的变换因子。在一元情形我们有 $f_Y(y) = f_X(h(y)) |h'(y)|$；二维情形类似，只是用行列式度量了面积缩放比例。

\section{2.3 条件分布：信息的价值}\label{sec:2.3}

\textbf{核心问题}：已知 $X_1$ 的值，$X_2$ 的分布会如何变化？

这就是条件分布要回答的问题。条件分布本质上是在 $X_1 = x_1$ 的条件下，重新审视 $X_2$ 的行为。

\textbf{离散型条件 pmf}：\footnote{推导见附录 \cref{sec:derivation-conditional-distribution}。}
\begin{equation}\label{eq:conditional-pmf-ch2}
p_{X_2 \mid X_1}(x_2 \mid x_1) = \frac{p_{X_1, X_2}(x_1, x_2)}{p_{X_1}(x_1)}, \quad p_{X_1}(x_1) > 0.
\end{equation}

\textbf{连续型条件 pdf}：
\begin{equation}\label{eq:conditional-pdf-ch2}
f_{X_2 \mid X_1}(x_2 \mid x_1) = \frac{f_{X_1, X_2}(x_1, x_2)}{f_{X_1}(x_1)}, \quad f_{X_1}(x_1) > 0.
\end{equation}

分子是联合密度，分母是边缘密度。这个比值衡量的是，在 $X_1 = x_1$ 的条件下，$(x_1, x_2)$ 这个"点"有多"典型"。

由此定义条件期望：
\begin{equation}\label{eq:conditional-expectation-ch2}
\mathbb{E}[u(X_2) \mid x_1] = \int_{-\infty}^{\infty} u(x_2) f_{X_2 \mid X_1}(x_2 \mid x_1) \, dx_2.
\end{equation}

\subsection{全期望公式与 Rao-Blackwell 定理}

这里有一个深刻的定理，它连接了条件与边缘：\footnote{推导见附录 \cref{sec:derivation-law-total-expectation}。}

\begin{equation}\label{th:law-total-expectation-ch2}
\mathbb{E}[\mathbb{E}(X_2 \mid X_1)] = \mathbb{E}(X_2).
\end{equation}

这称为\textbf{全期望公式}。它的含义是：\emph{对条件均值再取期望，等于无条件期望}。

更惊人的是方差的不等式：
\[
\text{var}[\mathbb{E}(X_2 \mid X_1)] \leq \text{var}(X_2).
\]\footnote{推导见附录 \cref{sec:derivation-variance-inequality}.}

\textbf{Rao-Blackwell 定理}：若用 $X_2$ 估计 $\mu_2 = \mathbb{E}X_2$，则条件均值 $E(X_2 \mid x_1)$ 比 $X_2$ 本身更精确（方差更小）。这在统计推断中极其重要——它告诉我们如何利用辅助信息改进估计。

\section{2.4 独立性：变量间的"无关系"}\label{sec:2.4}

\textbf{核心问题}：什么时候 $X_1$ 和 $X_2$ 互不影响？

若 $X_1$ 与 $X_2$ 独立，则知道 $X_1$ 的值不会改变我们对 $X_2$ 的认知。用条件分布的语言：
\[
f_{X_2 \mid X_1}(x_2 \mid x_1) = f_{X_2}(x_2) \quad \text{（不依赖于 } x_1 \text{）}.
\]

由此推出联合密度可分离：\footnote{推导见附录 \cref{sec:derivation-independence}。}

\begin{equation}\label{def:independence-ch2}
X_1 \text{ 与 } X_2 \text{ 独立} \iff f(x_1, x_2) \equiv f_1(x_1) f_2(x_2).
\end{equation}

一个实用的判据：

\begin{equation}\label{th:independence-factorization-ch2}
X_1 \text{ 与 } X_2 \text{ 独立} \iff f(x_1, x_2) \equiv g(x_1) h(x_2) \quad \text{（可分离变量）}.
\end{equation}

\textbf{独立性的后果}：
\begin{itemize}
\item $\mathbb{E}[g(X_1) h(X_2)] = \mathbb{E}[g(X_1)] \cdot \mathbb{E}[h(X_2)]$
\item $\text{cov}(X_1, X_2) = 0$
\end{itemize}

\begin remark}\label{rmk:cov-zero-not-independence}
\textbf{逆命题不成立！} $\text{cov}(X_1, X_2) = 0$ 并不意味着独立性（除二元正态分布外）。
\end remark}

这提醒我们：\textbf{不相关}（uncorrelated）和\textbf{独立}（independent）是两个不同概念。不相关只说明没有线性关系，但变量间可能存在非线性依赖。

\section{2.5 相关性：线性关联的度量}\label{sec:2.5}

我们已经知道 $\text{cov}(X_1, X_2) = 0$ 不能完全刻画独立性。那么如何量化 $X_1$ 与 $X_2$ 之间的"关联程度"？

协方差 $\text{cov}(X_1, X_2)$ 存在一个问题：它依赖于 $X_1, X_2$ 的量纲。例如，身高（cm）和体重（kg）的协方差与换算成英寸和磅时不同。

\textbf{解决方案}：先标准化，再求协方差：

\begin{equation}\label{def:correlation-ch2}
\rho_{12} = \text{corr}(X_1, X_2) = \frac{\text{cov}(X_1, X_2)}{\sqrt{\text{var}(X_1)} \sqrt{\text{var}(X_2)}} = \frac{\text{cov}(X_1, X_2)}{\sigma_1 \sigma_2}.
\end{equation}

这就是相关系数（correlation coefficient）。

\textbf{核心性质}：\footnote{推导见附录 \cref{sec:derivation-correlation-bounds}。}
\begin{itemize}
\item $-1 \leq \rho_{12} \leq 1$
\item $|\rho_{12}| = 1 \iff X_2 = aX_1 + b$（完美线性关系）
\item $\rho_{12} = 0$ 称\textbf{不相关}
\end{itemize}

\textbf{为什么是 $[-1, 1]$？} 这来自 Cauchy-Schwarz 不等式：
\[
|\text{cov}(X_1, X_2)| = |\mathbb{E}[(X_1-\mu_1)(X_2-\mu_2)]| \leq \sqrt{\mathbb{E}[(X_1-\mu_1)^2] \cdot \mathbb{E}[(X_2-\mu_2)^2]} = \sigma_1 \sigma_2.
\]

\begin remark}[相关 vs 协方差]
协方差取决于量纲（"米的协方差"vs"厘米的协方差"数值不同）；相关系数是无量纲的，消除了单位影响。因此，相关系数是更通用的"关联程度"度量。
\end remark}

\section{2.6 二元正态：最重要的二元分布}\label{sec:2.6}

在所有二元分布中，二元正态分布占据核心地位。它是一元正态分布在二维的自然推广。

\textbf{为什么重要？}
\begin{enumerate}
\item 边缘分布和条件分布都是正态的
\item 线性组合仍是正态
\item 不相关 $\iff$ 独立（这是正态独有的优雅性质）
\item 多元正态是计量经济学、信号处理等领域的基石
\end{enumerate}

\begin{definition}[二元正态分布]\label{def:bivariate-normal-ch2}
$(X_1, X_2)$ 服从二元正态分布，若其联合 pdf 为
\[
f(x_1, x_2) = \frac{1}{2\pi\sigma_1\sigma_2\sqrt{1-\rho^2}} \exp\left\{ -\frac{1}{2(1-\rho^2)} \left[ \frac{(x_1-\mu_1)^2}{\sigma_1^2} - 2\rho\frac{(x_1-\mu_1)(x_2-\mu_2)}{\sigma_1\sigma_2} + \frac{(x_2-\mu_2)^2}{\sigma_2^2} \right] \right\},
\]
记作 $(X_1, X_2) \sim N(\mu_1, \mu_2, \sigma_1^2, \sigma_2^2, \rho)$。
\end{definition}

\textbf{五个参数的含义}：
\begin{itemize}
\item $\mu_1, \mu_2$：边缘分布的均值
\item $\sigma_1^2, \sigma_2^2$：边缘分布的方差
\item $\rho$：相关系数
\end{itemize}

\begin remark}[条件分布的几何直觉]
在二元正态中，条件均值
\[
\mathbb{E}(X_2 \mid x_1) = \mu_2 + \rho\frac{\sigma_2}{\sigma_1}(x_1 - \mu_1)
\]
是 $x_1$ 的线性函数。这意味着：给定 $X_1 = x_1$，$X_2$ 的最佳预测也是线性的，其斜率由相关系数 $\rho$ 决定。
\end remark}

\textbf{关键性质}：

\begin{enumerate}
\item 边缘分布：$X_1 \sim N(\mu_1, \sigma_1^2)$，$X_2 \sim N(\mu_2, \sigma_2^2)$
\item 条件分布：
\[
X_2 \mid X_1 = x_1 \sim N\left(\mu_2 + \rho\frac{\sigma_2}{\sigma_1}(x_1 - \mu_1), \ \sigma_2^2(1-\rho^2)\right)
\]\footnote{推导见附录 \cref{sec:derivation-bivariate-normal-conditional}.}
\item $X_1 \Perp X_2 \iff \rho = 0$（正态独有的"不相关 $\Rightarrow$ 独立"）
\item 任意线性组合 $aX_1 + bX_2 \sim N(a\mu_1 + b\mu_2, \text{var}(aX_1 + bX_2))$
\end{enumerate}

\section{2.7 多元推广}\label{sec:2.7}

前述概念可以推广到 $n$ 个随机变量。记随机向量 $\mathbf{X} = (X_1, \ldots, X_n)'$，其联合 cdf 为
\[
F_{\mathbf{X}}(\mathbf{x}) = \mathbb{P}(X_1 \leq x_1, \ldots, X_n \leq x_n).
\]

\subsection{边缘与条件分布}

边缘 pdf 通过 $n-1$ 重积分得到：
\[
f_1(x_1) = \int_{-\infty}^{\infty} \cdots \int_{-\infty}^{\infty} f(x_1, x_2, \ldots, x_n) \, dx_2 \cdots dx_n.
\]

条件分布定义为联合密度与边缘密度之比：
\[
f_{2,\ldots,n|1}(x_2,\ldots,x_n \mid x_1) = \frac{f(x_1, x_2, \ldots, x_n)}{f_1(x_1)}.
\]

\subsection{多元独立性}

$n$ 个随机变量 $X_1, \ldots, X_n$ 相互独立，当且仅当
\[
f(x_1, x_2, \ldots, x_n) \equiv f_1(x_1) f_2(x_2) \cdots f_n(x_n).
\]

\textbf注意}：两两独立（pairwise independence）不一定推出相互独立（mutual independence）。Bernstein 的例子说明了这一点。

\subsection{线性组合的 mgf}

若 $X_1, \ldots, X_n$ 相互独立，$T = \sum_{i=1}^n k_i X_i$ 的 mgf 为
\[
M_T(t) = \prod_{i=1}^n M_i(k_i t).\footnote{推导见附录 \cref{sec:derivation-mgf-linear-combination}.}
\]

特别地，若 $X_1, \ldots, X_n$ iid（独立同分布），则 $T = \sum_{i=1}^n X_i$ 的 mgf 为
\[
M_T(t) = [M(t)]^n.
\]

\section{本章总结}\label{sec:chapter2-summary}

\begin{center}
\begin{tabular}{|c|c|}
\hline
\textbf{概念} & \textbf{公式/描述} \\ \hline
联合 cdf & $F_{X_1,X_2}(x_1,x_2) = \mathbb{P}(X_1\leq x_1, X_2\leq x_2)$ \\ \hline
边缘 pmf & $p_{X_1}(x_1) = \sum_{x_2} p_{X_1,X_2}(x_1,x_2)$ \\ \hline
边缘 pdf & $f_{X_1}(x_1) = \int_{-\infty}^{\infty} f(x_1,x_2)\,dx_2$ \\ \hline
条件 pmf & $p_{X_2|X_1}(x_2|x_1) = p(x_1,x_2)/p_{X_1}(x_1)$ \\ \hline
条件 pdf & $f_{X_2|X_1}(x_2|x_1) = f(x_1,x_2)/f_{X_1}(x_1)$ \\ \hline
全期望 & $\mathbb{E}[\mathbb{E}(X_2|X_1)] = \mathbb{E}X_2$ \\ \hline
独立性 & $f(x_1,x_2) = f_1(x_1)f_2(x_2)$ \\ \hline
协方差 & $\text{cov} = \mathbb{E}(X_1X_2) - \mathbb{E}X_1\mathbb{E}X_2$ \\ \hline
相关系数 & $\rho = \text{cov}/(\sigma_1\sigma_2)$，$-1 \leq \rho \leq 1$ \\ \hline
\end{tabular}
\end{center}

\section*{习题}\label{sec:exercises-ch2}

\begin{exercise}{2.1.1}
设 $X_1$ 和 $X_2$ 的联合 pdf 为 $f(x_1, x_2) = x_1 + x_2$，$0 < x_1 < 1$，$0 < x_2 < 1$，其他位置为零。求 $X_1$ 的边缘 pdf。
\end{exercise}

\begin{exercise}{2.2.1}
设 $(X_1, X_2)$ 的联合 pdf 为 $f(x_1, x_2) = 2$，$0 < x_1 < x_2 < 1$，其他位置为零。求 $Y = X_1 + X_2$ 的分布。
\end{exercise}

\begin{exercise}{2.3.1}
设 $X_1$ 和 $X_2$ 具有联合 pdf $f(x_1, x_2) = x_1 + x_2$，$0 < x_1 < 1$，$0 < x_2 < 1$。求给定 $X_1 = x_1$ 时 $X_2$ 的条件均值和条件方差。
\end{exercise}

\begin{exercise}{2.4.1}
证明：若 $X_1$ 和 $X_2$ 独立，则 $\mathbb{E}[g(X_1) h(X_2)] = \mathbb{E}[g(X_1)] \cdot \mathbb{E}[h(X_2)]$。
\end{exercise}

\begin{exercise}{2.5.1}
设 $(X, Y)$ 具有联合 pdf $f(x, y) = x + y$，$0 < x < 1$，$0 < y < 1$。求相关系数 $\rho$。
\end{exercise}

\begin{exercise}{2.6.1}
设 $X_1, X_2, X_3$ 为三个独立同分布的随机变量，每个的 pdf 为 $f(x) = 2x$，$0 < x < 1$。求 $Y = \max\{X_1, X_2, X_3\}$ 的分布。
\end{exercise}

% === 用户问答记录 ===%
\end{document}
